\subsection{The quasi-Galois extension generated by a set}

PROPOSITION 4. --- Let \((\mathbf{N}_i)_{i\in I}\) be a family of quasi-Galois extensions of K contained in \(\Omega\) . Let \(\mathbf{N} = \bigcap_{i\in I}\mathbf{N}_i\) and \(\mathbf{M} = \mathbf{K}\left(\bigcup_{i\in I}\mathbf{N}_i\right)\) ; then the extensions N and M

of K are quasi-Galois.

Let \(u\) be a K-automorphism of \(\Omega\) . We have \(u(\mathbf{N}_i) = \mathbf{N}_i\) for all \(i \in I\) , hence we obtain the equalities \(u(\mathbf{N}) = \mathbf{N}\) and \(u(\mathbf{M}) = \mathbf{M}\) ; now the proposition follows from Cor. 1 (V, p.~54).

Let A be a set of elements of \(\Omega\) , and let B be the set of elements of \(\Omega\) which are conjugates of elements of A; expressed differently, if G is the group of K-automorphisms of \(\Omega\) , we have \(\mathbf{B} = \bigcup_{u \in G} u(\mathbf{A})\) . For each \(u \in G\) we have

\(u(B) = B\) whence \(u(K(B)) = K(B)\) . Therefore \(K(B)\) is a quasi-Galois extension of K containing A and it is immediate that every quasi-Galois extension N of K containing A contains B, hence \(K(B)\) . We shall say that \(K(B)\) is the quasi-Galois extension generated by A. This definition applies in particular when A is an extension of K.

The next proposition follows directly from the preceding remarks.

PROPOSITION 5. --- Let E be an extension of K contained in Ω and N the quasi-Galois extension generated by E. If A is a subset of Ω such that E = K(A) then N = K(B), where B is the set of elements of Ω which are conjugate to an element of A.

COROLLARY 1. --- If E is an extension of finite degree of K, then the quasi-Galois extension N of K generated by E is of finite degree over K.

For we have \(E = K(A)\) where A is finite, hence the set B of conjugates of elements of A is finite, whence the corollary follows by Th. 2 (V, p.~18).

COROLLARY 2. --- Every quasi-Galois extension N of K is the union of quasi-Galois subextensions of N of finite degree over K.

Let \(x \in \mathbb{N}\) and let \(\mathbf{N}_x\) be the quasi-Galois extension of \(\mathbf{K}\) generated by \(\{x\}\) . Since \(\mathbf{K}(x)\) is of finite degree over \(\mathbf{K}\) , the same is true of \(\mathbf{N}_x\) (Cor. 1) and we have \(x \in \mathbf{N}_x\) .

